\section{Hierarchies of Group Languages Are Not Collapsed by \(\rext\)}\label{sec:group}

The main result of this section is Corollary~\ref{cor:strict-nilp-solv},
which is a consequence of Theorem~\ref{thm:group-robust} below and states that
the classes of regular languages approximable by nilpotent/solvable group languages form strict hierarchies.
This is completely different from the situation in Theorem~\ref{thm:sf},
which shows that the dot-depth hierarchy is collapsed by \(\rext\).

\begin{theorem}[Y.]\label{thm:group-robust}
  If \(\CC \subseteq \group_A\) is a local variety of group languages, then \(\rext(\CC) \cap \group_A = \CC\).
  In particular, if \(\CC \subsetneq \CD \subseteq \group_A\) are local subvarieties,
  then \(\rext(\CC) \cap \CD = \CC\), \(\CD \not\subseteq \rext(\CC)\), and \(\rext(\CC) \subsetneq \rext(\CD)\).
\end{theorem}
Our previous result \cite[Theorem 1]{gcommeasure} is the special case of Theorem~\ref{thm:group-robust} where \(\CC = \gcom_A\).

\begin{proof}
  Let \(L \in \group_A \setminus \CC\).
  It suffices to show that \(L\) is \(\CC\)-immeasurable.
  Since \(\CC\) is a local variety and \(L \notin \CC\), Theorem~\ref{thm:duality} implies
  that there exists an equation \(u = v\) (where \(u, v \in \widehat{A^*}\)) that every language in \(\CC\) satisfies but \(L\) does not.
  Let \(G\) be the syntactic monoid of \(L\), which is a group, and \(\eta\colon A^* \to G\) be the syntactic morphism of \(L\).
  By Lemma~\ref{lem:group-profinite-identity},
  there exists \(w \in \widehat{A^*}\) such that the equation \(w = \varepsilon\) is equivalent to \(u = v\) on any \(A\)-generated finite group.
  Hence, for every \(L_1 \in \CC\) and its syntactic morphism \(\theta_1\colon A^* \to G_1\),
  we have \(\hat{\theta}_1(w) = \hat{\theta}_1(\varepsilon) = 1_{G_1}\) and \(\hat{\eta}(w) \neq \hat{\eta}(\varepsilon) = 1_G\).
  Since \(\hat{\eta}(w) \neq 1_G\) and \(G\) is the syntactic monoid of \(L\), there exist \(x, y \in A^*\)
  such that
  \begin{equation}
    \eta(x) \hat{\eta}(w) \eta(y) \in \eta(L) \centernot\iff \eta(x y) \in \eta(L). \label{eq:ebc45b27}
  \end{equation}
  Since the two languages \(\eta^{-1}(\eta(x)), \eta^{-1}(\eta(y)) \subseteq A^*\) are group languages,
  they are non-null by Lemma~\ref{lem:groupdensity}.
  Hence the concatenation \(I \defeq \eta^{-1}(\eta(x)) \eta^{-1}(\eta(y))\) is also non-null (by Lemma~\ref{lem:density}, for example).
  To complete the proof, it suffices to show that \(L_1, L_2 \in \CC\) and \(L_1 \subseteq L \subseteq L_2\) implies \(I \subseteq L_2 \setminus L_1\).
  Let \(\theta_1\colon A^* \to G_1\) and \(\theta_2\colon A^* \to G_2\) be the syntactic morphism of \(L_1\) and \(L_2\), respectively.
  Let \((w_n)_{n \in \nat}\) be a Cauchy sequence in \(A^*\) converging to \(w\) in \(\widehat{A^*}\).
  Then we have \(\eta(w_k) = \hat{\eta}(w)\), \(\theta_1(w_k) = \hat{\theta}_1(w)\), and \(\theta_2(w_k) = \hat{\theta}_2(w)\) for sufficiently large \(k \in \nat\).

  We show \(I \subseteq L_2 \setminus L_1\).
  Let \(z \in I\).
  Then there exist \(\tilde{x} \in \eta^{-1}(\eta(x))\) and \(\tilde{y} \in \eta^{-1}(\eta(y))\) such that \(z = \tilde{x} \tilde{y}\).
  Let \(i \in \{1, 2\}\).
  Since \(\theta_i(w_k) = 1_{G_i}\) and \(G_i\) is the syntactic monoid of \(L_i\),
  we have \(\theta_i(\tilde{x} w_k \tilde{y}) \in \theta_i(L_i) \iff \theta_i(\tilde{x} \tilde{y}) \in \theta_i(L_i)\),
  \ie,
  \begin{equation}
    z = \tilde{x} \tilde{y} \in L_i \iff \tilde{x} w_k \tilde{y} \in L_i. \label{eq:fbaddcb1}
  \end{equation}
  It follows from \eqref{eq:ebc45b27}, \(\eta(\tilde{x}) = \eta(x)\), \(\eta(\tilde{y}) = \eta(y)\), and \(\hat{\eta}(w) = \eta(w_k)\) that
  \begin{equation}
    z = \tilde{x} \tilde{y} \in L \centernot\iff \tilde{x} w_k \tilde{y} \in L. \label{eq:500d85c6}
  \end{equation}
  %It is easy to show that \eqref{eq:fbaddcb1}, \eqref{eq:500d85c6},
  %and \(L_1 \subseteq L \subseteq L_2\) implies \(z \in L_2 \setminus L_1\).
  We show \(z \in L_2 \setminus L_1\) by considering two cases depending on whether \(z \in L\) or \(z \notin L\).
  First, suppose that \(z \in L\).
  Since \(L \subseteq L_2\), we have \(z \in L_2\).
  On the other hand, by \eqref{eq:500d85c6}, we know that \(\tilde{x} w_k \tilde{y} \notin L \supseteq L_1\), which implies \(\tilde{x} w_k \tilde{y} \notin L_1\).
  Then, by \eqref{eq:fbaddcb1}, it follows that \(z \notin L_1\), and thus we obtain \(z \in L_2 \setminus L_1\).
  The case where \(z \notin L\) can be proved in a similar manner.
\end{proof}

\begin{corollary}\label{cor:strict-nilp-solv}
  Suppose \(\card{A} \geq 2\).
  For \(n \geq 1\), we have
  \begin{itemize}
    \item \(\rext(\gnil_A^{\leq n}) \subsetneq \rext(\gnil_A^{\leq n + 1}) \subsetneq \rext(\group_A)\), and
    \item \(\rext(\gsol_A^{\leq n}) \subsetneq \rext(\gsol_A^{\leq n + 1}) \subsetneq \rext(\group_A)\).
  \end{itemize}
\end{corollary}
\begin{proof}
  By Proposition~\ref{prop:strict} and Theorem~\ref{thm:group-robust}.
\end{proof}

